% Options for packages loaded elsewhere
\PassOptionsToPackage{unicode}{hyperref}
\PassOptionsToPackage{hyphens}{url}
\PassOptionsToPackage{dvipsnames,svgnames,x11names}{xcolor}
%
\documentclass[
  10pt,
  a4paper,
]{article}
\usepackage{amsmath,amssymb}
\usepackage{iftex}
\ifPDFTeX
  \usepackage[T1]{fontenc}
  \usepackage[utf8]{inputenc}
  \usepackage{textcomp} % provide euro and other symbols
\else % if luatex or xetex
  \usepackage{unicode-math} % this also loads fontspec
  \defaultfontfeatures{Scale=MatchLowercase}
  \defaultfontfeatures[\rmfamily]{Ligatures=TeX,Scale=1}
\fi
\usepackage{lmodern}
\ifPDFTeX\else
  % xetex/luatex font selection
\fi
% Use upquote if available, for straight quotes in verbatim environments
\IfFileExists{upquote.sty}{\usepackage{upquote}}{}
\IfFileExists{microtype.sty}{% use microtype if available
  \usepackage[]{microtype}
  \UseMicrotypeSet[protrusion]{basicmath} % disable protrusion for tt fonts
}{}
\makeatletter
\@ifundefined{KOMAClassName}{% if non-KOMA class
  \IfFileExists{parskip.sty}{%
    \usepackage{parskip}
  }{% else
    \setlength{\parindent}{0pt}
    \setlength{\parskip}{6pt plus 2pt minus 1pt}}
}{% if KOMA class
  \KOMAoptions{parskip=half}}
\makeatother
\usepackage{xcolor}
\usepackage[margin=1in]{geometry}
\usepackage{color}
\usepackage{fancyvrb}
\newcommand{\VerbBar}{|}
\newcommand{\VERB}{\Verb[commandchars=\\\{\}]}
\DefineVerbatimEnvironment{Highlighting}{Verbatim}{commandchars=\\\{\}}
% Add ',fontsize=\small' for more characters per line
\newenvironment{Shaded}{}{}
\newcommand{\AlertTok}[1]{\textcolor[rgb]{1.00,0.00,0.00}{\textbf{#1}}}
\newcommand{\AnnotationTok}[1]{\textcolor[rgb]{0.38,0.63,0.69}{\textbf{\textit{#1}}}}
\newcommand{\AttributeTok}[1]{\textcolor[rgb]{0.49,0.56,0.16}{#1}}
\newcommand{\BaseNTok}[1]{\textcolor[rgb]{0.25,0.63,0.44}{#1}}
\newcommand{\BuiltInTok}[1]{\textcolor[rgb]{0.00,0.50,0.00}{#1}}
\newcommand{\CharTok}[1]{\textcolor[rgb]{0.25,0.44,0.63}{#1}}
\newcommand{\CommentTok}[1]{\textcolor[rgb]{0.38,0.63,0.69}{\textit{#1}}}
\newcommand{\CommentVarTok}[1]{\textcolor[rgb]{0.38,0.63,0.69}{\textbf{\textit{#1}}}}
\newcommand{\ConstantTok}[1]{\textcolor[rgb]{0.53,0.00,0.00}{#1}}
\newcommand{\ControlFlowTok}[1]{\textcolor[rgb]{0.00,0.44,0.13}{\textbf{#1}}}
\newcommand{\DataTypeTok}[1]{\textcolor[rgb]{0.56,0.13,0.00}{#1}}
\newcommand{\DecValTok}[1]{\textcolor[rgb]{0.25,0.63,0.44}{#1}}
\newcommand{\DocumentationTok}[1]{\textcolor[rgb]{0.73,0.13,0.13}{\textit{#1}}}
\newcommand{\ErrorTok}[1]{\textcolor[rgb]{1.00,0.00,0.00}{\textbf{#1}}}
\newcommand{\ExtensionTok}[1]{#1}
\newcommand{\FloatTok}[1]{\textcolor[rgb]{0.25,0.63,0.44}{#1}}
\newcommand{\FunctionTok}[1]{\textcolor[rgb]{0.02,0.16,0.49}{#1}}
\newcommand{\ImportTok}[1]{\textcolor[rgb]{0.00,0.50,0.00}{\textbf{#1}}}
\newcommand{\InformationTok}[1]{\textcolor[rgb]{0.38,0.63,0.69}{\textbf{\textit{#1}}}}
\newcommand{\KeywordTok}[1]{\textcolor[rgb]{0.00,0.44,0.13}{\textbf{#1}}}
\newcommand{\NormalTok}[1]{#1}
\newcommand{\OperatorTok}[1]{\textcolor[rgb]{0.40,0.40,0.40}{#1}}
\newcommand{\OtherTok}[1]{\textcolor[rgb]{0.00,0.44,0.13}{#1}}
\newcommand{\PreprocessorTok}[1]{\textcolor[rgb]{0.74,0.48,0.00}{#1}}
\newcommand{\RegionMarkerTok}[1]{#1}
\newcommand{\SpecialCharTok}[1]{\textcolor[rgb]{0.25,0.44,0.63}{#1}}
\newcommand{\SpecialStringTok}[1]{\textcolor[rgb]{0.73,0.40,0.53}{#1}}
\newcommand{\StringTok}[1]{\textcolor[rgb]{0.25,0.44,0.63}{#1}}
\newcommand{\VariableTok}[1]{\textcolor[rgb]{0.10,0.09,0.49}{#1}}
\newcommand{\VerbatimStringTok}[1]{\textcolor[rgb]{0.25,0.44,0.63}{#1}}
\newcommand{\WarningTok}[1]{\textcolor[rgb]{0.38,0.63,0.69}{\textbf{\textit{#1}}}}
\setlength{\emergencystretch}{3em} % prevent overfull lines
\providecommand{\tightlist}{%
  \setlength{\itemsep}{0pt}\setlength{\parskip}{0pt}}
\setcounter{secnumdepth}{-\maxdimen} % remove section numbering
\ifLuaTeX
  \usepackage{selnolig}  % disable illegal ligatures
\fi
\usepackage{bookmark}
\IfFileExists{xurl.sty}{\usepackage{xurl}}{} % add URL line breaks if available
\urlstyle{same}
\hypersetup{
  colorlinks=true,
  linkcolor={blue},
  filecolor={Maroon},
  citecolor={Blue},
  urlcolor={blue},
  pdfcreator={LaTeX via pandoc}}

\author{}
\date{}

\begin{document}

\section{BA0 Base Analysis working hypothesis -
R2}\label{ba0-base-analysis-working-hypothesis---r2}

\textbf{Status: STRONG CANDIDATE / NOT CLOSED}\\
\textbf{Supersedes for active BA0 reasoning:}
\texttt{BA0\_BASE\_ANALYSIS\_WORKING\_HYPOTHESIS\_R1.md}\\
\textbf{Evidence incorporated:} BA0-R closure, BA0-T1, BA0-T2.

\subsection{1. Working responsibility
statement}\label{1-working-responsibility-statement}

For a governed documentation baseline, Base Analysis provides an
\textbf{accepted, methodology-neutral analyzable representation of
shared project meaning} grounded in governed documentation and explicit
reviewed analytical additions.

Its responsibility is to make that meaning reusable across:

\begin{itemize}
\tightlist
\item
  progressive human understanding and source drill-down;
\item
  multiple analysis consumers and method-specific projections;
\item
  explicit provenance and source-localized diagnostics;
\item
  change impact, revalidation and targeted re-analysis;
\item
  governed corrective feedback from analysis back to project
  documentation.
\end{itemize}

Base Analysis is not project authority. Governed documentation is.

\subsection{2. Canonicality clarified}\label{2-canonicality-clarified}

\texttt{Canonical} means accepted semantic identity and meaning
\textbf{within the analysis layer for a governed baseline}. It does not
require:

\begin{itemize}
\tightlist
\item
  one permanently stored graph;
\item
  one graph database;
\item
  one serialization;
\item
  one diagram;
\item
  one modeling language implementation.
\end{itemize}

A Base Analysis may be reused or reproducibly rebuilt, provided accepted
identities, provenance, review decisions and source alignment remain
stable enough for the intended analyses and change comparisons.

\subsection{3. Authority boundary}\label{3-authority-boundary}

\begin{Shaded}
\begin{Highlighting}[]
\NormalTok{governed documentation  = project authority}
\NormalTok{Base Analysis            = accepted analytical representation}
\NormalTok{projection/view          = derived presentation or method input}
\NormalTok{analysis interpretation  = method{-}owned semantics}
\NormalTok{candidate finding/gap    = reviewable analysis output}
\NormalTok{corrective candidate     = proposed governed{-}document change}
\end{Highlighting}
\end{Shaded}

No derived layer may silently replace the authority of governed
documentation.

\subsection{4. Working invariants}\label{4-working-invariants}

\subsubsection{BA0-I1 - Source
authority}\label{ba0-i1---source-authority}

Accepted Base Analysis claims that state project facts must be grounded
in governed documentation or explicitly identified as reviewed
analytical additions. Analysis output cannot silently create governed
project truth.

\subsubsection{BA0-I2 - Analysis-layer semantic
canonicality}\label{ba0-i2---analysis-layer-semantic-canonicality}

For a governed baseline, accepted analytical referents and propositions
must have stable enough identity and meaning to be reused across views,
consumers and re-analysis. Physical persistence is an implementation
choice.

\subsubsection{BA0-I3 - Explicit origin and
provenance}\label{ba0-i3---explicit-origin-and-provenance}

Grounded statements, derivations, reviewed additions and unresolved
diagnostics must retain origin and source traceability sufficient for
review and drill-down.

\subsubsection{BA0-I4 - No silent inference
authority}\label{ba0-i4---no-silent-inference-authority}

NLP, LLMs, similarity, graph construction, tool defaults or analyst
shortcuts may propose candidate structure but cannot silently establish
identity, equivalence, ownership, project facts or governed commitments.

\subsubsection{BA0-I5 - Method
neutrality}\label{ba0-i5---method-neutrality}

The common Base Analysis core contains shared project semantics, not
STRIDE, STRIDE-AI or other method-owned threat categories. Methods may
add isolated interpretation in their own analysis layer.

\subsubsection{BA0-I6 - Representation
independence}\label{ba0-i6---representation-independence}

The responsibility statement is independent from graph databases,
diagrams, DFD notation, SysML, JSON/YAML schemas or one renderer.
Concrete representations must preserve the accepted semantics they claim
to realize.

\subsubsection{BA0-I7 - Bounded modeling
scope}\label{ba0-i7---bounded-modeling-scope}

Base Analysis includes only shared semantics justified by project
understanding, analysis reuse, provenance, diagnostics or
change/re-analysis needs. It is not a general-purpose digital twin or
complete systems model by default.

\subsubsection{BA0-I8 - Change alignment and scoped
revalidation}\label{ba0-i8---change-alignment-and-scoped-revalidation}

When governed documentation changes, affected analytical semantics and
dependent analyses must be identifiable for review or re-execution
without automatically invalidating unrelated work.

\subsubsection{BA0-I9 - Diagnostic visibility and
localization}\label{ba0-i9---diagnostic-visibility-and-localization}

Contradiction, ambiguity, missing context or insufficient analysis basis
must not be silently canonicalized. A diagnostic should preserve enough
source/referent/context information to localize targeted review.

\subsubsection{BA0-I10 - No method/tool semantic
ownership}\label{ba0-i10---no-methodtool-semantic-ownership}

ThreatForge, a renderer, an extraction pipeline, a graph engine or a
methodology plugin may implement/support DDTA contracts but does not
establish the DDTA semantic core by implementation accident.

\subsubsection{BA0-I11 - Progressive
teachability}\label{ba0-i11---progressive-teachability}

Accepted Base Analysis semantics must support non-authoritative
overview-to-detail projections that let a human progressively understand
the project and reach exact governed sources without rereading the whole
corpus first.

This is a representational responsibility, not a claim that any specific
visualization has empirically proven lower reading time.

\subsubsection{BA0-I12 - Governed corrective
feedback}\label{ba0-i12---governed-corrective-feedback}

Analysis may emit source-targeted corrective documentation candidates
for gaps, contradictions, missing decisions, missing requirements or
other justified issues. Only governed review may accept the change as
project truth; accepted source changes trigger Base Analysis
revalidation.

\subsubsection{BA0-I13 - Consumer isolation with shared
meaning}\label{ba0-i13---consumer-isolation-with-shared-meaning}

A specific analysis consumer may select/project the Base Analysis facts
it needs and attach method-owned interpretation, but it should not have
to reconstruct the same project meaning independently or redefine shared
semantics in order to operate.

\subsection{5. Working origin classes}\label{5-working-origin-classes}

The exact material schema remains open, but BA0 still needs to
distinguish at least:

\begin{itemize}
\tightlist
\item
  \textbf{grounded} - directly supported by governed documentation;
\item
  \textbf{derived} - analytical structure/statement justified from
  grounded sources without a new project commitment;
\item
  \textbf{reviewed analytical addition} - analysis-layer addition
  accepted for analytical use but not pretending to be governed project
  truth;
\item
  \textbf{diagnostic / unresolved} - conflict, ambiguity, missing
  information or insufficiency that prevents a justified canonical
  choice.
\end{itemize}

A diagnostic is not intended to become a permanent substitute for
repairing governed documentation when the source itself needs
correction.

\subsection{6. Progressive human
understanding}\label{6-progressive-human-understanding}

Human-facing views are projections of accepted Base Analysis semantics.
A useful project navigation path is:

\begin{Shaded}
\begin{Highlighting}[]
\NormalTok{project / macro concerns}
\NormalTok{        {-}\textgreater{} decisions and responsibility boundaries}
\NormalTok{        {-}\textgreater{} functional behavior}
\NormalTok{        {-}\textgreater{} specialized constraints}
\NormalTok{        {-}\textgreater{} method{-}relevant view}
\NormalTok{        {-}\textgreater{} exact proposition / source clause}
\end{Highlighting}
\end{Shaded}

Different views may omit different information, but omission must not
create contradictory authored models.

\subsection{7. Specific-analysis
consumption}\label{7-specific-analysis-consumption}

A methodology-specific consumer should be able to:

\begin{enumerate}
\def\labelenumi{\arabic{enumi}.}
\tightlist
\item
  declare/select the shared semantics it requires;
\item
  build a method-specific projection or interpretation;
\item
  keep its taxonomy/evidence payload outside the common core;
\item
  return results, insufficiencies or documentation gaps with provenance
  to the shared semantics and governed sources.
\end{enumerate}

BA0-T2 demonstrates this boundary only with a bounded STRIDE transfer
probe. Universal methodology support is not claimed.

\subsection{8. Governed feedback loop}\label{8-governed-feedback-loop}

\begin{Shaded}
\begin{Highlighting}[]
\NormalTok{governed documentation}
\NormalTok{ {-}\textgreater{} Base Analysis}
\NormalTok{ {-}\textgreater{} human/method projection}
\NormalTok{ {-}\textgreater{} analysis result or diagnostic}
\NormalTok{ {-}\textgreater{} corrective documentation candidate}
\NormalTok{ {-}\textgreater{} governed review}
\NormalTok{ {-}\textgreater{} accepted correction / clarification / supersession}
\NormalTok{ {-}\textgreater{} Base Analysis revalidation or rebuild}
\NormalTok{ {-}\textgreater{} impacted analysis re{-}execution}
\end{Highlighting}
\end{Shaded}

Not every analysis result requires a documentation change, and not every
documentation gap must become a SecurityRequirement. The semantic owner
and document kind depend on the issue discovered. The thesis may
separately evaluate a narrower accepted-Finding-to-SecurityRequirement
path.

\subsection{9. Explicit non-goals}\label{9-explicit-non-goals}

BA0 is not responsible for:

\begin{itemize}
\tightlist
\item
  reliable automatic extraction from arbitrary narrative/legacy
  documentation;
\item
  silently accepting LLM/NLP outputs;
\item
  defining the final BAE type taxonomy;
\item
  copying a complete existing systems-modeling language;
\item
  becoming a STRIDE DFD or another threat-method model by definition;
\item
  proving universal support for all analysis methods;
\item
  enforcing one persisted graph/storage technology;
\item
  encoding every useful semantic responsibility as a first-class
  metaclass;
\item
  empirically proving that one visualization improves human
  comprehension or effort.
\end{itemize}

\subsection{10. Open questions reserved for later
phases}\label{10-open-questions-reserved-for-later-phases}

Still open after T2:

\begin{itemize}
\tightlist
\item
  which semantic responsibilities become first-class BAE types versus
  relations/properties/projections;
\item
  exact identity/equivalence and lifecycle rules;
\item
  minimal relation vocabulary;
\item
  material provenance schema and source locators;
\item
  formal consumer/overlay contract;
\item
  exact view types and projection contracts;
\item
  broader method-neutrality and holdout regression.
\end{itemize}

These belong to BA1-BA6 only after BA0 responsibility closes.

\subsection{11. Status and next step}\label{11-status-and-next-step}

R2 is a \textbf{STRONG CANDIDATE}, not BA0 closure.

The next authorized microstep is
\texttt{BA0-T3\ -\ responsibility\ and\ non-goals\ closure\ review}. BA1
must not start until BA0 is explicitly closed.

\end{document}
